\documentclass[11pt]{article}

\usepackage[T1]{fontenc}
\usepackage[utf8]{inputenc}
\usepackage[spanish]{babel}
\usepackage{graphicx}
\usepackage{enumitem}

\parindent = 0cm
\setlist{itemsep=0.6em}

\usepackage{amsmath}
\usepackage{amssymb,amsfonts,latexsym,cancel}
\providecommand{\abs}[1]{\lvert#1\rvert}
\providecommand{\norm}[1]{\LVert#1\RVert}
\usepackage{multirow}
\usepackage[table,xcdraw]{xcolor}

% Márgenes limpios y compactos para todo el documento
\usepackage[lmargin=2cm,rmargin=2cm,top=2cm,bottom=2.5cm]{geometry}
\usepackage{xspace}

% ======================================================================
%  DEFINICIÓN DEL TÍTULO DEL TEMA
% ======================================================================
\newcommand{\titulotema}{Cuentas Nacionales y Mediciones Macroeconómicas II}

% ======================================================================
%  CONFIGURACIÓN DE ENCABEZADOS Y PIES DE PÁGINA (fancyhdr)
% ======================================================================
\usepackage{fancyhdr}

% --- Estilo general: desde la página 2 en adelante ---
\pagestyle{fancy}
\fancyhf{}
\fancyhead[L]{Macroeconomía (580323)}
\fancyhead[R]{\titulotema}
\renewcommand{\headrulewidth}{0.5pt} % Sin línea superior
\fancyfoot[C]{\thepage}
\fancyfoot[R]{Semestre 2026-2}
\fancyfoot[L]{CMP/IGU/AFR}
\renewcommand{\footrulewidth}{0.5pt}

% --- Estilo exclusivo de la Página 1 ---
\fancypagestyle{primerapagina}{%
  \fancyhf{}
  \fancyhead[L]{\small FACULTAD DE INGENIERÍA\\DEPTO DE INGENIERÍA INDUSTRIAL \\ MACROECONOMÍA (580323)}
  \fancyhead[R]{\includegraphics[scale=0.13]{escudoudec.jpg}}
  \renewcommand{\headrulewidth}{0.9pt}
  \fancyfoot[C]{\thepage}
  \fancyfoot[R]{Semestre 2026-2}
  \fancyfoot[L]{CMP/IGU/AFR}
  \renewcommand{\footrulewidth}{0.5pt}
}

% Enlaces
\usepackage{url}
\usepackage{hyperref}
\hypersetup{
    colorlinks=true,
    linkcolor=blue,
    filecolor=magenta,
    urlcolor=cyan,
}

% Paquete de recuadros avanzados
\usepackage[many]{tcolorbox}

% ======================================================================
%  DEFINICIÓN DEL RECUADRO CON BORDES REDONDEADOS Y BARRA AZUL
% ======================================================================
\definecolor{headblue}{RGB}{0, 15, 137}      
\definecolor{frameblue}{RGB}{41, 98, 153}     
\definecolor{bgsoftblue}{RGB}{247, 249, 253}  

\newtcolorbox{solucion}[1][]{%
    enhanced,
    breakable,
    colback=bgsoftblue,
    colframe=frameblue!75!black,
    colbacktitle=headblue,
    coltitle=white,
    fonttitle=\bfseries\normalsize,
    title={Solución:\ifx&#1&\else\ \normalfont\small(#1)\fi},
    arc=2.5mm,
    boxrule=0.7pt,
    titlerule=0pt,
    toptitle=1.5mm,
    bottomtitle=1.5mm,
    left=4mm, right=4mm, top=3mm, bottom=3mm,
    before skip=4mm,
    after skip=4mm,
    skin first=enhanced,
    skin middle=enhanced,
    skin last=enhanced,
    arc=2.5mm,
    bottomrule=0.7pt,
    toprule=0.7pt
}

\newenvironment{solucionpartes}{%
    \begin{enumerate}[label=\textbf{(\alph*)}, leftmargin=1.2cm, itemsep=0.5em]
}{%
    \end{enumerate}
}

\newcommand{\checkformula}[1]{\noindent\textbf{Chequeo de consistencia: } #1 \quad \checkmark}

% ======================================================================
%  COMANDOS PERSONALIZADOS
% ======================================================================

\newcounter{ejercicio}
\newcommand{\ejercicio}[1][]{%
  \stepcounter{ejercicio}%
  \bigskip\noindent\textbf{Ejercicio \arabic{ejercicio}\ifx&#1&\relax.\else\ (#1).\fi}\ %
}

\newenvironment{partes}
  {\begin{enumerate}[label=(\alph*), leftmargin=1.5cm]}
  {\end{enumerate}}

\newenvironment{tabladatos}
  {\begin{center}}
  {\end{center}}

\newenvironment{recuadro}[1]{%
  \noindent\colorbox{black!8}{%
    \parbox{\dimexpr\linewidth-2\fboxsep\relax}{%
      \centering\Large\bfseries #1
    }%
  }\par\vspace{4mm}\noindent
}{\par}

\begin{document}
\thispagestyle{primerapagina}

\vspace*{0.05cm}

\begin{center}
  {\Large \textbf{Listado 2 -- Macroeconomía (580323)}}\\[0.15cm]
  {\large \titulotema}
\end{center}

\bigskip

\textbf{Profesor:} Cristian Mardones, Ph.D. \quad
\textbf{Ayudantes:} Ignacio Garrido U y Alex Fierro R.

\bigskip

%=======================================================================
% EJERCICIO 1 (Certamen 2013) -- Cuentas Nacionales
%=======================================================================
\ejercicio[Certamen 2013]
Los datos siguientes corresponden a información de cuentas nacionales (en miles):

\begin{tabladatos}
$\text{PIB} = 60{,}000,\quad \text{Inversión Bruta} = 12{,}000,\quad \text{Inversión Neta} = 9{,}000,$\\[0.2em]
$\text{Consumo} = 33{,}000,\quad \text{Gasto del Gobierno} = 14{,}500,\quad \text{Superávit Fiscal} = 500.$
\end{tabladatos}

\noindent
Calcule el \textbf{Producto Interno Neto}, las \textbf{Exportaciones Netas}, los \textbf{impuestos menos transferencias}, el \textbf{ingreso personal disponible} y el \textbf{ahorro}.



%=======================================================================
% EJERCICIO 2 (Examen 2021) -- Contabilidad Nacional: pares de transacciones
%=======================================================================
\ejercicio[P, Examen 2021]
Mencione cuál es la diferencia desde el punto de vista de la contabilidad nacional para las dos transacciones económicas de cada enunciado:

\begin{partes}
  \item Una empresa que compra combustible para que sus vendedores realicen reparto en camioneta y una empresa que le paga a sus vendedores un bono de monto equivalente para que ellos mismos compren el combustible.
  \item Usted decide comprar un pisco peruano en vez de un pisco chileno.
  \item Un hospital público que compra una jeringa y una familia que compra ese mismo tipo de jeringa en una farmacia.
\end{partes}



%=======================================================================
% EJERCICIO 3 (Certamen 2013) -- Identidad Ahorro-Inversión y Déficit Fiscal
%=======================================================================
\ejercicio[P, Certamen 2013]
A partir de la definición de producto igual a la demanda agregada, demuestre la relación entre ahorro externo, ahorro del gobierno y ahorro privado. Si el gobierno incurre en un déficit fiscal y el ahorro privado es menor que la inversión, ¿qué implicancia tiene esto sobre las exportaciones netas?



%=======================================================================
% EJERCICIO 4 (Certamen 2024) -- Modelo Macroeconómico de Producción y Cuentas
%=======================================================================
\ejercicio[Certamen 2024]
Suponga que la economía de un país se puede representar por las siguientes ecuaciones:
\begin{tabladatos}
$Y = 15K^{0,5}L^{0,5}, \quad K = 2500, \quad L = 400$\\[0.4em]
$C = 250 + 0,6Y_D, \quad T = 0,1Y, \quad TR = 1000, \quad G = 1600, \quad XN = 400$
\end{tabladatos}

\noindent
Determine:

\begin{partes}
  \item PIB, Ingreso Disponible, Ahorro Privado, Ahorro del Gobierno e Inversión.
\end{partes}



\newpage




%=======================================================================
% EJERCICIO 6 -- PIB vs. PNB y críticas a la medición del PIB (teórico)
%=======================================================================
\ejercicio[P]
El Banco Central de un país anuncia con orgullo que el PIB creció un 4\% este año. Un periodista de El Mercurio (cuerpo EyN) le pide a usted, como asesor económico, que comente la noticia.

\begin{partes}
  \item Explique la diferencia conceptual entre el \textbf{Producto Interno Bruto (PIB)} y el \textbf{Producto Nacional Bruto (PNB)}. Dé un ejemplo de una actividad productiva en Chile en la cual ambas medidas puedan diferir de forma relevante.
  \item Un estudio para Chile (Mardones \& del Río, 2019) muestra que el PIB se reduciría en 11,3\% si se descontaran los daños ambientales asociados a la minería del cobre. ¿Cuál de las tres principales críticas a la medición del PIB revisadas en clases ilustra mejor este resultado? Explique brevemente en qué consiste dicha crítica.
  \item Nombre y explique brevemente \textbf{otra} de las críticas a la medición del PIB vistas en clases (distinta a la de la parte anterior), dando un ejemplo distinto al de la minería.
  \item Si el PIB de un país crece 4\% pero, al mismo tiempo, aumenta fuertemente el trabajo doméstico no remunerado y el empleo informal no registrado, ¿el bienestar de la población necesariamente creció en la misma magnitud que el PIB? Justifique.
\end{partes}



%=======================================================================
% EJERCICIO 7 -- PIB real: año base fijo vs. índices de volumen encadenados
%=======================================================================
\ejercicio[P]
Una economía produce solamente dos bienes: computadores y bicicletas. Las cantidades producidas y sus precios entre 2024 y 2026 se muestran en la tabla:

\begin{tabladatos}
\begin{tabular}{lcccccc}
\hline\hline
\multirow{2}{*}{\textbf{Bien}} & \multicolumn{2}{c}{\textbf{2024}} & \multicolumn{2}{c}{\textbf{2025}} & \multicolumn{2}{c}{\textbf{2026}} \\
\cline{2-3} \cline{4-5} \cline{6-7}
 & Precio & Cantidad & Precio & Cantidad & Precio & Cantidad \\
\hline
Computadores (unid.) & \$1.000 & 50 & \$900 & 70 & \$850 & 90 \\
Bicicletas (unid.) & \$200 & 300 & \$220 & 280 & \$240 & 260 \\
\hline\hline
\end{tabular}
\end{tabladatos}

\begin{partes}
  \item Calcule el PIB nominal para 2024, 2025 y 2026.
  \item Calcule el PIB real de cada año utilizando \textbf{2024 como año base fijo}, y la tasa de crecimiento real 2024--2025 y 2025--2026 con este método.
  \item Calcule ahora el PIB real mediante un \textbf{índice de volumen encadenado}: para cada tramo, valorice el crecimiento usando los precios del año inicial de ese mismo tramo (2024 para el tramo 2024--2025, y 2025 para el tramo 2025--2026), y luego encadene ambas tasas de crecimiento a partir del PIB nominal de 2024.
  \item Compare el nivel del PIB real de 2026 obtenido en (b) y en (c). ¿Por qué difieren, si ambos parten del mismo año base? Relacione su respuesta con el cambio en los precios relativos de computadores y bicicletas entre 2024 y 2026.
  \item ¿Por qué la literatura recomienda actualmente el uso de índices de volumen encadenados por sobre un año base fijo y lejano en el tiempo?
\end{partes}


\newpage
%=======================================================================
% EJERCICIO 8 -- Comercio exterior y tipo de cambio: Viña El Boldo
%=======================================================================
\ejercicio
La Viña El Boldo es la única productora de vino de una pequeña economía abierta. Durante el año produjo 400.000 litros de vino, vendidos en el mercado interno a un precio de \$2.800 por litro. De esa producción, 250.000 litros se destinan al consumo de los hogares del país, y los 150.000 litros restantes se exportan a un precio internacional de US\$4,20 por litro.

Para elaborar el vino, la viña importa barricas de roble francés como insumo del proceso productivo: 500 barricas a US\$300 cada una. Además, la economía importa 20 tractores agrícolas a US\$25.000 cada uno, los cuales se destinan íntegramente a inversión del sector. El tipo de cambio vigente es de \$950 CLP por dólar.

Exprese todos sus resultados en pesos chilenos (CLP) y determine:

\begin{partes}
  \item El \textbf{valor bruto de la producción (VBP)} de vino.
  \item El valor de las importaciones de \textbf{insumos} (barricas) y, descontando dicho valor, el \textbf{Producto Interno Bruto (PIB)} generado por la viña.
  \item El valor del \textbf{consumo} de los hogares y el valor de las \textbf{exportaciones}.
  \item El valor de la \textbf{inversión} asociada a la importación de tractores y las \textbf{exportaciones netas (XN)} del período.
  \item Suponga que, en lugar de dejar flotar su moneda, el Banco Central de este país decide fijar el tipo de cambio en \$900 CLP por dólar y comprometerse a comprar y vender divisas a ese precio sin límite. ¿Qué régimen cambiario estaría adoptando el país? Mencione un ejemplo real de un país que haya adoptado un esquema similar (o uno aún más extremo) y explique brevemente en qué consiste.
\end{partes}




%=======================================================================
% EJERCICIO 9 (Certamen 2025) -- Contabilidad del Crecimiento y Productividad Total de Factores
%=======================================================================
\ejercicio[P, Certamen 2025]
Suponga que la producción final de la economía se puede caracterizar por la siguiente función de producción:

\[
Y_t = A_t K_t^{0,6} L_t^{0,4}
\]

\noindent
A partir de la ecuación previa, obtenga paso a paso una fórmula para determinar la \textbf{tasa de crecimiento de la productividad total de factores}. Luego, use los siguientes datos para calcular la evolución de la productividad total de factores a través del tiempo:

\begin{tabladatos}
\begin{tabular}{cccc}
\hline\hline
\textbf{Período} & \textbf{Crec. PIB (\%)} & \textbf{Crec. Capital (\%)} & \textbf{Crec. Trabajo (\%)} \\
\hline
2001--2005 & 3,50 & 2,00 & 1,00 \\
2006--2010 & 2,80 & 1,80 & 0,80 \\
2011--2015 & 2,50 & 2,00 & 0,60 \\
2016--2020 & 2,40 & 2,20 & 0,50 \\
2021--2025 & 2,20 & 2,00 & 1,00 \\
\hline\hline
\end{tabular}
\end{tabladatos}



%=======================================================================
% EJERCICIO 10 -- Del crecimiento del PIB al bienestar por persona
%  (continuación creativa del Ejercicio 9: PTF + PIB per cápita)
%=======================================================================
\ejercicio
El PIB total puede crecer por más razones que el simple aumento de la productividad: también puede crecer porque hay más trabajadores, o porque la población crece más rápido o más lento que el empleo. Para separar estos efectos, continuemos con la economía del \textbf{Ejercicio 8}, cuya función de producción es
\[
Y_t = A_t K_t^{0,6} L_t^{0,4}.
\]

\noindent
Defina el \textbf{capital por trabajador} $k_t = K_t/L_t$ y la \textbf{razón empleo-población} $e_t = L_t/N_t$, donde $N_t$ es la población total del país. Denote la tasa de crecimiento porcentual de una variable $X$ como $\hat{X}$ (por ejemplo, $\hat{Y}$ es el crecimiento porcentual del PIB).

\begin{partes}
  \item Demuestre que el \textbf{PIB per cápita}, $y_t = Y_t/N_t$, puede escribirse como
  \[
  y_t = A_t\, k_t^{0,6}\, e_t.
  \]
  \emph{Ayuda:} reemplace $K_t = k_t L_t$ dentro de la función de producción antes de dividir por $N_t$.

  \item A partir del resultado anterior, demuestre que la tasa de crecimiento del PIB per cápita satisface
  \[
  \hat{y} \;=\; \hat{A} \;+\; 0{,}6\cdot\hat{k} \;+\; \hat{e}.
  \]
  Interprete económicamente cada uno de los tres términos del lado derecho: ¿qué tipo de fenómeno económico representa cada uno (piense en productividad, ahorro/inversión, y demografía)?

  \item Usando los valores de la tabla que ya calculó en el \textbf{Ejercicio 8}, y los siguientes datos adicionales de crecimiento poblacional, complete para cada período los valores de $\hat{k}$ (crecimiento del capital por trabajador), $\hat{e}$ (crecimiento de la razón empleo-población) y $\hat{y}$ (crecimiento del PIB per cápita):

  \begin{tabladatos}
  \begin{tabular}{ccccc}
  \hline\hline
  \textbf{Período} & \textbf{Crec. Población (\%)} & \textbf{$\hat{k}$} & \textbf{$\hat{e}$} & \textbf{$\hat{y}$} \\
  \hline
  2001--2005 & 1,30 & & & \\
  2006--2010 & 1,10 & & & \\
  2011--2015 & 1,00 & & & \\
  2016--2020 & 0,90 & & & \\
  2021--2025 & 0,60 & & & \\
  \hline\hline
  \end{tabular}
  \end{tabladatos}

  \item Verifique que sus resultados de $\hat{y}$ en (c) coinciden con calcular directamente $\hat{y} = \hat{Y} - \hat{N}$ usando los datos de crecimiento del PIB del Ejercicio 8. ¿Por qué esta igualdad \textbf{debe} cumplirse siempre, independientemente de los valores numéricos?

  \item ¿En cuál de los cinco períodos el crecimiento del PIB per cápita fue más ``sano'', en el sentido de estar impulsado principalmente por la productividad (PTF) y no por la acumulación de capital ni por factores demográficos? Justifique con los datos de su tabla.

  \item Suponga que, manteniendo todo lo demás constante, la tasa de crecimiento poblacional de 2021--2025 hubiese sido de 1,30\% en vez de 0,60\% (por ejemplo, por mayor inmigración neta sin que esos nuevos habitantes se incorporaran de inmediato a la fuerza laboral). Sin recalcular todo, explique qué le ocurriría a $\hat{e}$ y a $\hat{y}$ de ese período, y por qué.
\end{partes}

\end{document}